% !TEX root = ../notes.tex

\documentclass[../notes.tex]{subfiles}

\begin{document}

\section{March 18}
Here we go.

\subsection{More on Strata}
Last time, we explained how to take a Weyl conjugacy class and produce an element of $\op{Conj}U_*$, which of course index the strata. 
\begin{remark}
	If we restrict to the ``elliptic'' Weyl elements, then the conjugacy classes inject into the strata. We may call the image of this map the ``elliptic'' strata.
\end{remark}
\begin{remark}
	If $G$ is semisimple, then for any $w$ of minimal length in an elliptic conjugacy class, the length is equal to the dimension of the centralizer of the corresponding stratum. Professor Lusztig has only checked this in a case-by-case manner.
\end{remark}
\begin{remark}
	Except for one stratum in type $E_8$, each elliptic stratum contains a semisimple element.
\end{remark}
\begin{example}
	In type $A_n$, the unique elliptic conjugacy class is given by the $n$-cycle. It goes to the regular semisimple stratum.
\end{example}
Let's discuss the proof of the following conjecture of Lusztig.
\begin{theorem}[Carnovale]
	Let $G$ be a reductive group. Any stratum of $G$ is locally closed.
\end{theorem}
\begin{proof}
	Let $G_m$ be the collection of $g\in G$ of dimension $m$. This is a union of strata, and it is known to be locally closed. %For example, one can realize $G_m$ as the image of some algebraic map; more explicitly, $G_m$ should be a family of conjugacy classes.

	The proof will involved ``pieces.'' Let's review this definition.
	\begin{notation}
		Let $G$ be a reductive group. For each $g\in G$, let $g_s$ be the semisimple element, and we define $H_G(g)\coloneqq Z_G(g_s)^\circ$.
	\end{notation}
	\begin{remark}
		Note that $H_G(g)$ is the Levi subgroup of some parabolic because it is the centralizer of a torus.
	\end{remark}
	\begin{definition}[piece]
		Let $G$ be a reductive group. For any Levi subgroup $L$ of a parabolic, we may let $S_1$ be the pre-image of an isolated class in $L/Z^\circ(L)$. Here, an element of a connected reductive group is said to be isolated if and only if the connected components of the centralizer of its semisimple part is not contained in a Levi subgroup of a proper parabolic subgroup. We can then define $Y_{L,S_1}$ to be the union of all conjugacy classes in $G$ which meet those $g\in S_1$ and having $H_G(g)=L$. It turns out that there are only finitely many pairs $(L,S_1)$ up to conjugacy, and these $Y_{L,S_1}$s are the pieces of $G$.
	\end{definition}
	\begin{remark}
		Lusztig has shown that the closure of a piece is a union of pieces.
	\end{remark}
	Now, Carnovale shows that a stratum $\Sigma$ is a union of pieces. In fact, for any piece $J\subseteq\Sigma$, one can find $m$ with $\Sigma\subseteq G_m$, and then one can show that $\overline J\cap G_m\subseteq\Sigma$. It follows that $\Sigma$ is the union of these closed subsets $\overline J\cap G_m$, so $\Sigma$ is closed in $G_m$ and hence locally closed in $G$.
\end{proof}
\begin{remark}
	In fact, the proof has shown that each stratum $\Sigma$ is closed in $G_m$. One can show further that a stratum is a union of irreducible components.
\end{remark}
% Recall that we discussed how to send irreducible representations of $W$ to $\op{Conj}U_p$, which then embeds into $\op{Irr}W$; we denote its image by $\op{Irr}_p(W)$.
\begin{remark}
	There is a map $\varepsilon_p$ from irreducible representations to $W$ to $\op{Conj}U_p$: any irreducible representation of $W$ is parameterized by the Springer correspondence by a unipotent element and a local system, so we can forget about the local system. We can then return another irreducible representation by adding back in the trivial local system. For an irreducible representation $E$ of $W$, we may write $\varepsilon_*(E)$ for the value of $\varepsilon_p(E)$ if this irreducible representation is independent of $p$; if it is not independent of $p$, then it only changes at a single prime $p_0$, and we define $\varepsilon_*(E)$ to be this exceptional value.
\end{remark}
\begin{remark}
	There is a notion of a ``superstratum.'' One can decompose $\op{Irr}W$ into sets $\op{Irr}_cW$, where $c$ is a two-sided cell. For example, in type $C$, one has the symbols $(0,1,2)$ and $(0,2,1)$ and $(1,0,2)$ in the same two-sided cell. A superstratum is then the union of the strata $\Sigma$ which fix a two-sided cell, meaning that there is a map $\delta\colon G\to\op{Irr}_*W$ (maybe given by the Springer correspondence), and the superstratum is the union of the $\Sigma$ for which $\delta(\Sigma)\in c$, for fixed two-sided cell $c$. It is conjectured that each superstratum is locally closed and that its irreducible components are rationally smooth.
\end{remark}
\begin{example}
	In type $A$, the strata agree with the superstrata.
\end{example}
\begin{example}
	For $\op{Sp}_4$, the superstrata are given as follows.
	\begin{itemize}
		\item There is a superstratum of the zero-dimensional conjugacy classes.
		\item There is a superstratum of the eight-dimensional conjugacy classes.
		\item Everything else belongs to a third superstratum.
	\end{itemize}
\end{example}

\end{document}